\section{Serialization to \texttt{.qrd}}\label{sec:serialization}

The on-disk \texttt{.qrd} format must carry the same shift the rest of this design
makes everywhere else: it stops recording resolved geometry and starts recording
\emph{intent}. Today a child's spatial relationship to its parent is persisted as an
opaque \code{<offset x y z/>} triple --- the serialized form of the CM-to-CM
\cppinline|Vector3| stored in \code{childParts}. That triple inherits every pathology
discussed in \cref{sec:datamodel}: it is a resolved coordinate, it conflates a part's center
of mass with its placement, and it bakes a particular set of part lengths
into the file so that a later geometry edit silently invalidates it. Replacing the
storage type (\cref{sec:datamodel}) without replacing its serialized projection would
merely move the bug onto disk. This section defines the replacement element, the
reader/writer changes, and why the format version need not gate the loader.

\subsection{The \texttt{<link>} element, kept inline}

The \code{<offset>} child element is replaced, per child, by an inline \code{<link>}
element that serializes the \seat{}-bearing \cppinline|StationLink| rather than a
resolved vector. The element is attached to the child \code{<part>} it governs, exactly
where \code{<offset>} sat, and it carries the four authored quantities of the link:
the parent station, the child station, the gap, and the seat kind.

\begin{designrule}[The link is stored inline with its child, not in a sibling block]
A tempting alternative --- collecting all parent$\rightarrow$child relationships into a
single sibling \code{<joints>} block keyed by part name --- is a hazard, because part
\emph{names need not be unique}: \code{getName()} is explicitly documented as
non-identifying (\src{model/parts/Part.h}{193}; identity is the per-instance
\code{getId()}, which is not serialized). A name-keyed joints block would therefore be
ambiguous the moment a design contains two parts with the same name, which is a
permitted and common state. Anchoring each \code{<link>} to its owning child element
sidesteps name resolution entirely: the ownership tree in the XML \emph{is} the keying.
\end{designrule}

\Cref{lst:qrd} shows the target \texttt{0.2} encoding of the \code{xl75\_multi} design,
with each of its three non-root attachments expressed as a \code{<link>}.

\begin{listing}[htbp]
\begin{minted}{xml}
<QtRocketDesign version="0.2">         <!-- minor bump; loader still accepts 0.1 -->
  <part type="NoseCone" name="MultiNose">
    <params baseRadius="0.0395" length="0.30" wallThickness="0"
            density="1700" solid="true"/>
    <!-- root: no <link> -->
    <children>
      <part type="BodyTube" name="MultiBody">
        <params innerRadius="0.0376" outerRadius="0.0395" length="0.90"
                density="1700"/>
        <link seat="Abut" parentStation="0.0" childStation="1.0" gap="0"/>
        <!-- ^ equals the default link; omittable -->
        <children>
          <part type="BodyTube" name="MultiCoupler">
            <params innerRadius="0.036" outerRadius="0.0376" length="0.08"
                    density="1700"/>
            <link seat="NestInBore" parentStation="1.0" childStation="0.0"
                  gap="0.04"/>
          </part>
          <part type="FinSet" name="MultiFins">
            <params rootChord="0.10" tipChord="0.04" span="0.06" sweep="0.04"
                    thickness="0.003" bodyRadius="0.0395" finCount="6"
                    density="1700"/>
            <link seat="OnSurface" parentStation="0.06" childStation="0.0"
                  gap="0"/>
          </part>
        </children>
      </part>
    </children>
  </part>
  <motor commonName="M1350W"/>
  <sim .../>
</QtRocketDesign>
\end{minted}
\caption{The \code{xl75\_multi} design in the target \texttt{0.2} format. Each child's
spatial relationship is an intent-bearing \code{<link>}, never a resolved coordinate.
The root carries no link; the body's link equals the default and could be omitted.}
\label{lst:qrd}
\end{listing}

\subsection{Writer: \texttt{writeOffset} becomes \texttt{writeLink}}

The current writer emits the resolved offset through \code{writeOffset}
(\src{model/DesignSerializer.cpp}{53}), which puts \code{x}/\code{y}/\code{z} attributes
from a \cppinline|Vector3|; \code{writePart} adds that node as the child's \code{offset}
sub-element. Under the new model this helper is replaced by
\code{writeLink(const StationLink\&)}, which emits the four link attributes ---
\code{seat}, \code{parentStation}, \code{childStation}, and \code{gap} --- in place of
the coordinate triple. The \seat{} enumerator is mapped to and from its attribute string
through a single shared lookup table, so the writer and the reader agree on the
spelling (\seat{Abut}, \seat{NestInBore}, \seat{OnSurface}) by construction rather than
by two independently maintained switch statements.

\begin{keyinsight}[A default-equal link is omittable]
Because the zero-config default is a meaningful physical relationship --- the child's
fore plane abutting the parent's aft plane (\cref{sec:authoring}) --- a link whose
fields all equal that default carries no information and may be omitted from the file
entirely. The writer can elide it, and the reader (below) reconstructs the same default.
This keeps the common nose$\rightarrow$body$\rightarrow$\ldots stack as terse on disk as
it is in the authoring API, and it makes the serialized form proportional to the design's
actual complexity rather than to its part count.
\end{keyinsight}

\subsection{Reader: \texttt{buildPart} consumes the link}

The reader's recursion lives in \code{buildPart} (\src{model/DesignSerializer.cpp}{117}).
Today it reconstructs a \cppinline|Vector3| from the child's
\code{offset.<xmlattr>.x/y/z} (defaulting each to \code{0.0}) and calls the legacy
\code{addChildPart(child, Vector3)}. Under the new model it instead reads the link
attributes --- \code{parentStation}, \code{childStation}, \code{seat}, \code{gap} ---
defaulting them respectively to \code{0.0}, \code{1.0}, \seat{Abut}, and \code{0},
then constructs a \cppinline|StationLink{...}| and calls
\cppinline|addChildPart(child, StationLink{...})|. Those defaults are exactly the
zero-config link, so a minimal \code{<part>} carrying \emph{no} \code{<link>} at all
still attaches by abutting its parent --- the omittable-link property of the writer and
the defaulting behavior of the reader are two halves of the same contract.

The reader inherits the existing fail-closed discipline. A \code{seat} string that maps
to no \seat{} enumerator is a load error, raised on the same path as today's
\code{makePart} rejection of an unknown part type, so a malformed or hand-edited file is
refused with a clear diagnostic rather than silently coerced to \seat{Abut}.

\begin{designrule}[Resolved coordinates are never serialized]
The file stores only authored intent --- station fractions, a gap, and a seat kind. The
absolute \code{Pose} of each part is recomputed by the resolver (\cref{sec:resolver}) on
load. This is the on-disk expression of the ``derive, never store'' rule: a \texttt{.qrd}
file can be edited (lengthen a tube, swap a cone) and reloaded, and the seams move to
follow the new geometry because nothing in the file pins them to the old. It is also the
natural validation point for migration --- recomputing the placement on load is precisely
when the overlap sweep (\cref{sec:diagnostics}) gets to run against a freshly read tree.
\end{designrule}

\subsection{Version: writer emits \texttt{0.2}; no loader gate change}

The format version moves from \texttt{0.1} to \texttt{0.2} to mark the
\code{<offset>}$\rightarrow$\code{<link>} change. This is a minor bump because the
loader's acceptance rule is already keyed on the \emph{major} component only: the gate
(\src{model/DesignSerializer.cpp}{186}) splits the version on the dot and rejects only an
unrecognized major, so it admits any minor within major~\texttt{0}. Both \texttt{0.1}
(legacy \code{<offset>}) and \texttt{0.2} (\code{<link>}) therefore pass the same gate
with no change to it; the reader distinguishes the two element shapes structurally, by
which element is present on the child, not by inspecting the version string.

\begin{hardfail}[The current writer still emits \texttt{0.1}]
In the present working tree the writer hard-codes the version string \texttt{"0.1"}
(\src{model/DesignSerializer.cpp}{156}), and \code{writeOffset} still emits the
coordinate triple. The \texttt{0.2}/\code{<link>} encoding described above is the
\emph{target} of this design, not the current behavior; the only file-format work the
loader gate needs is none, precisely because \src{model/DesignSerializer.cpp}{186}
already tolerates the minor bump.
\end{hardfail}
